\section{A* Search}

A* is a shortest-path search algorithm that extends Dijkstra's algorithm
by incorporating an estimate of the remaining distance to the goal.

For each node $n$, A* considers three quantities:

\[
g(n),
\]

the known cost of the shortest path discovered so far from the start node
to $n$,

\[
h(n),
\]

a heuristic estimate of the remaining cost from $n$ to the goal, and

\[
f(n)=g(n)+h(n),
\]

the estimated total cost of a path from the start to the goal that passes
through $n$.

While Dijkstra's algorithm selects the next node using only $g(n)$, A*
selects the node with the smallest value of $f(n)$.

This allows A* to use information about the location of the goal when
deciding which part of the graph to explore.

\subsection{Euclidean-Distance Heuristic}

The visibility graph is embedded in two-dimensional Euclidean space.
Therefore, a natural heuristic is the straight-line distance from a node
to the goal.

For a node

\[
n=(x_n,y_n)
\]

and goal

\[
G=(x_G,y_G),
\]

the heuristic is

\[
h(n)
=
\sqrt{
(x_G-x_n)^2
+
(y_G-y_n)^2
}.
\]

This gives the shortest possible geometric distance between the node and
the goal without considering obstacles.

\subsection{Admissibility}

For A* to guarantee an optimal path, the heuristic should not overestimate
the true remaining shortest-path cost.

A heuristic satisfying

\[
h(n)\leq h^*(n)
\]

is called \emph{admissible}, where $h^*(n)$ denotes the true shortest-path
cost from node $n$ to the goal.

For the visibility graph, every edge weight is a Euclidean distance.
Any path from $n$ to the goal therefore consists of one or more Euclidean
line segments.

By the triangle inequality,

\[
\|n-G\|
\leq
\|n-v_1\|
+
\|v_1-v_2\|
+
\cdots
+
\|v_k-G\|.
\]

Therefore, the straight-line Euclidean distance to the goal can never
exceed the length of a valid path through the visibility graph.

Hence,

\[
h(n)=\|n-G\|
\]

is an admissible heuristic.

\subsection{Consistency}

A stronger property than admissibility is \emph{consistency}.

A heuristic is consistent if, for every graph edge between nodes $u$ and
$v$,

\[
h(u)
\leq
w(u,v)+h(v).
\]

For the Euclidean visibility graph,

\[
h(u)=\|u-G\|,
\]

\[
w(u,v)=\|u-v\|,
\]

and

\[
h(v)=\|v-G\|.
\]

The triangle inequality gives

\[
\|u-G\|
\leq
\|u-v\|+\|v-G\|.
\]

Therefore, the Euclidean-distance heuristic is also consistent.

With a consistent heuristic, once a node is selected for expansion using
the smallest $f$-value, its best path cost does not need to be reconsidered
under the standard A* formulation used in this project.

\subsection{Relaxation in A*}

Although A* selects nodes using

\[
f(n)=g(n)+h(n),
\]

edge relaxation still operates on the actual path cost $g(n)$.

For an edge from $u$ to $v$, the candidate path cost is

\[
g_{\text{candidate}}
=
g(u)+w(u,v).
\]

If

\[
g_{\text{candidate}}<g(v),
\]

then

\[
g(v)=g_{\text{candidate}},
\]

and the parent of $v$ is updated to $u$.

The heuristic does not change the actual cost of a path. It only influences
which node A* chooses to explore next.

\subsection{Relationship to Dijkstra's Algorithm}

Dijkstra's algorithm can be viewed as a special case of A* in which

\[
h(n)=0
\]

for every node.

In that case,

\[
f(n)
=
g(n)+0
=
g(n),
\]

so A* selects nodes using exactly the same criterion as Dijkstra's
algorithm.

Therefore,

\[
\boxed{
\text{Dijkstra}
=
\text{A* with } h(n)=0
}
\]

This relationship makes the two algorithms particularly useful to compare
within the same visibility graph.